\section{Introduction}
\label{sec:intro}

In February 2022 SpaceX launched 49 Starlink satellites into a low parking orbit. Right after launch a
minor geomagnetic storm arrived, the air at that altitude became denser, and 38 of the
satellites were dragged back into the atmosphere before they could climb
away~\cite{fang2022starlink}. Two years later, the May 2024 ``Gannon'' storm, the strongest in
about twenty years, sped up the decay of almost every tracked object in low Earth orbit and set
off a wave of manoeuvres by operators~\cite{parker2024gannon}. Drag decides how long each of
the tens of thousands of objects in LEO stays up, storms can change it within hours, and
collision screening and re-entry forecasts depend on getting it right.

The physics is well known. During a storm, energy from the solar wind is dumped into the upper
atmosphere at high latitudes. The thermosphere heats and expands, the density at a given height
goes up, and so does drag. Putting a number on the storm's effect is harder than it sounds,
though, because the Sun also changes the atmosphere in a slower way. Its ultraviolet (UV)
output, tracked by the F10.7 radio flux, heats the thermosphere too, and periods with high UV
also tend to have more storms. If we simply correlate storm strength with decay, part of the UV
effect gets credited to the storms. On top of that, operators predict density with empirical
models such as NRLMSIS~\cite{picone2002nrlmsise,emmert2021nrlmsis,emmert2022msis21}, so it
matters just as much whether these models are wrong during storms.

Our approach is to treat each storm as a natural experiment. Nobody can schedule when a storm
arrives, so if we adjust for the measured F10.7 and the season, the decay just before a storm
and the decay just after it should differ only because of the storm. This is a causal reading,
and it rests on one assumption: that no time-varying confounder is left after the adjustment.
Our placebo and control checks support that assumption but cannot prove it. We use only passive
objects (debris and rocket bodies) in the main analysis, because they cannot fire thrusters and
their orbits change only through natural forces. We also measure their decay against a physics
baseline that already includes the real F10.7, so much of the UV effect is removed before any
statistics are run.

We ask three questions.
\textbf{Q1:} How much extra decay do storms cause, and how wrong is NRLMSIS~2.1 during them?
\textbf{Q2:} Which objects are hit hardest, by altitude, inclination and drag sensitivity?
\textbf{Q3:} How much of the storm's altitude loss do operators undo, during the storm and
after it?

To answer them we built a 250-object passive cohort covering 20 storms and estimated the storm
effect two ways, with a stacked event study and with double machine learning, which agree. We
estimated the NRLMSIS~2.1 storm-time error separately and checked the three largest storms
with a calm-to-calm comparison that does not depend on how TLE fits spread a loss across days.
A causal forest turned the per-object effects into a vulnerability map, and for Starlink we
built an exploratory offset ratio that splits the operator response into a part during the
storm and a part afterwards. Every estimator was tested on a simulator with a planted effect.

Sections~\ref{sec:related} to~\ref{sec:design} cover related work, data, the physics baseline
and the causal design, Sections~\ref{sec:q1} to~\ref{sec:q3} answer the three questions,
Section~\ref{sec:robust} tries to break the results and Section~\ref{sec:discussion} discusses
their limits.
